The two baselines lost $F_1$ in different ways. On Internal (Figure~\ref{fig:exp3_pr_scatter}), BLINKER favoured recall over precision, with a mean recall of 95.18\% and a mean precision of 67.94\%. On Cao2018 (Figure~\ref{fig:exp3_pr_scatter}), the same bias was stronger, with a mean recall of 97.15\% against a mean precision of 56.96\%. MNE, in contrast, showed no systematic bias, since its mean recall and precision differed by less than 2 percentage points on both datasets. Its operating points were instead widely dispersed, and recall fell below 30\% in several sessions on each dataset. The proposed pipeline had higher mean precision and recall than MNE on both datasets. Its mean recall was lower than BLINKER's (90.14\% on Internal and 89.09\% on Cao2018), but its mean precision was higher (88.03\% and 75.91\%). The recall bias of BLINKER is consistent with the pooled-threshold mechanism described above, and the low-recall MNE sessions are consistent with a range-based threshold raised by large excursions elsewhere in the recording. Neither mechanism was isolated, however. The proposed pipeline sets its threshold at the median plus three robust standard deviations of the flagged epochs, so BLINKER's threshold differs from it in centre and multiplier as well as in signal pool. The thresholds estimated by the three methods were also not recorded.
